% Methods. Draft prose --- target ~800 words; verify symbols at revision.

\paragraph{Detectors.}
Given a question $q$, the target model $M$ samples $N$ answers
$\{y_1,\dots,y_N\}$ at temperature $T$. Bidirectional natural-language
inference partitions them into semantic-equivalence clusters
$\mathcal{C}(q)$: $y_i$ and $y_j$ share a cluster iff each entails the other
under the NLI model. Semantic entropy is the Shannon entropy of the cluster-size
distribution,
\[
H_{\mathrm{SE}}(q) = -\sum_{c \in \mathcal{C}(q)} p(c)\log p(c),
\qquad p(c) = \frac{|c|}{N}.
\]
Semantic Reformulation Entropy (SRE)~\citep{tong2025sre} first generates
$R$ paraphrases of $q$, draws $K$ samples from each, pools all $R{\times}K$
answers, and computes a single entropy over the joint clustering. A higher score
is read as higher uncertainty; a detector flags an answer as a likely
hallucination when its score exceeds a threshold.

\paragraph{Threat model.}
The attacker is black-box: it may query $M$ and the detector but has no
gradients. It may replace the question $q$ with any paraphrase $q'$ subject to a
hard equivalence constraint, $\mathrm{NLI}(q,q')$ and $\mathrm{NLI}(q',q)$ both
entailment, so $q'$ asks the same thing and admits the same answer. The model's
answer is not otherwise constrained; the attack manipulates only the
\emph{detector's score}.

\paragraph{Attack objectives.}
We define two directions. The \emph{false-alarm} attack targets a question the
model answers correctly and seeks to \emph{inflate} the detector score so a
correct answer is flagged:
$\max_{q'} H(q')$ subject to the equivalence constraint.
The \emph{hide} attack targets a question the model answers incorrectly and seeks
to \emph{suppress} the score so a hallucination is missed:
$\min_{q'} H(q')$, i.e. $\max_{q'} -H(q')$. Writing both as maximization of a
signed objective lets a single optimizer serve both, with the sign encoding the
direction. The same objective is instantiated against either detector
($H \in \{H_{\mathrm{SE}}, H_{\mathrm{SRE}}\}$), which is what makes the
cross-detector transfer measurement well-defined.

\paragraph{Optimizer.}
We adapt the constraint-preserving zeroth-order search of
SECA~\citep{liang2025seca}. The search maintains a beam of the top-$P$ paraphrase
candidates; each round proposes $M$ further paraphrases per beam member via an
LLM paraphraser, scores them with the objective, keeps those that improve on the
running best, and admits a candidate only if it passes the bidirectional-NLI gate
against the \emph{original} $q$ (not the parent, so meaning cannot drift across
rounds). The search terminates on a score threshold or an iteration budget. We
swap SECA's objective (confidence in a wrong multiple-choice token) for the
detector entropy above, and its LLM-judge feasibility check for the NLI gate.

\paragraph{Target selection.}
Both attacks require a pool of questions with a known clean outcome---wrong
answers for the hide attack, correct answers for the false-alarm attack. We label
correctness by greedy decoding under an alias-aware span-match oracle (canonical
normalization, crediting any accepted answer form) and draw a
\emph{score-independent}, stratified-random sample within each correctness
stratum. Selection reads only the correctness label and a fixed seed, never any
detector score; consequently the SE and SRE campaigns operate on the
\emph{identical} question set by construction, which is what makes the
cross-detector comparison well-defined. This is a deliberate departure from
selecting targets at extreme clean entropy: extreme selection inflates the clean
detector AUROC toward $1.0$ by construction---the selection rule reflected
back---and would overstate any measured degradation. We therefore report
degradation against the clean AUROC of the \emph{fair} pool, not the extreme one.

\paragraph{Success criterion.}
An attack on a question succeeds only if all three conditions hold: (i) the
paraphrase $q'$ passes the bidirectional-NLI equivalence gate; (ii) the detector
score moves in the intended direction by at least $\delta$ nats past its clean
value; and (iii) the model's hallucination \emph{status} is unchanged under
$q'$---still wrong for the hide attack, still correct for the false-alarm
attack---checked with the same greedy decode and span oracle used to label the
pool. Condition (iii) is essential: a score can fall because the paraphrase
accidentally elicited a \emph{correct} answer (removing the hallucination rather
than hiding it), or rise because it elicited a wrong one; without the status
check these would be miscounted as attacks. We report the status-gated success
rate alongside the weaker entropy-only rate so the attrition between them is
explicit.

\paragraph{Evaluation and reporting.}
We report every rate with a bootstrap $95\%$ confidence interval over questions.
Beyond the aggregate score shift we report operating-point flips: fixing a
detector threshold at a target false-positive rate on clean data, we count hide
targets that move from flagged to unflagged and false-alarm targets that move
from unflagged to flagged---the operationally meaningful events. We additionally
break out the subset of would-be successes voided by condition~(iii) because the
model's answer flipped, since these are candidates for equivalence-gate leakage
and bound the fidelity of the NLI constraint.

\paragraph{Transfer and defense.}
To test whether the vulnerability is paradigm-level rather than implementation-
specific, we take a paraphrase optimized against one detector and measure the
score change it induces on another (SE$\leftrightarrow$SRE), and, as a stretch,
on a hidden-state probe (SEP~\citep{kossen2024seps}) that never samples. As a
first defense we average the detector score over $k$ paraphrases of whatever
input is received, on the intuition that a single adversarial phrasing is diluted
by its own paraphrases; this mirrors SRE's built-in reformulation and lets us ask
whether that mechanism confers attack resistance.
